\section*{Aufgabe 2a)}

Der Bearbeitung der Aufgabe 2 liegt die Funktion \texttt{homo\_ocs.m} zu Grunde,
welche den Hamilton-Operator mit Potential in Abhängigkeit von α in Matrix-Darstellung
zurück gibt. Sie ist in \lref{2funk} abgedruck.

\lstinputlisting[label=lst:2funk,caption={homo\_ocs.m}]{../code/homo_osc.m}

Zunächst sollten für verschiedene α die drei niedrigsten Eigenwerte bestimmt werden.
Der Quellcode ist hierfür in \lref{2a} zu finden.
\lstinputlisting[firstline=64,lastline=70,label=lst:2a,caption={steuerung.m}]{../code/steuerung.m}

\begin{lstlisting}[caption=Ausgabe von \lref{2a},label=2aa]
eigen_values =
    alpha=1   alpha=2   alpha=3   alpha=4
    0.5087    0.5000    0.5115    0.5302
    1.1691    1.5000    1.7253    1.8998
    1.6239    2.5000    3.1851    3.7278
\end{lstlisting}

In der Ausgabe \lref{2aa} sind die Ergebnisse dargestellt. Die 3 niedrigsten Eigenwerte
für ein α befinden sich jeweils in einer Spalte. Man kann sehr gut erkennen, dass man für
$α = 2$ das analytische Ergebniss erhält. Ein Vergleich mit den Daten aus CP1 ist
weiterhin nicht sinnvoll, da das Potential anders definiert war. Für den Spezialfall
$α = 2$ deckt sich zufällig das damals definierte Potential mit dem in für diese
Übung verwendeten. Das Ergebnis für die Eigenwerte lautete damals:
\begin{equation*}
0.500000,\;1.500000,\;2.499999
\end{equation*}
Beim Vergleich stellen wir als fest, dass die in dieser Übung verwendete Methode
zur Bestimmung der Eigenwerte genauer ist. Damals wurde eine Differentialgleichung
mit dem Runge-Kutta-Verfahren gelöst.



\section*{Aufgabe 2b)}

Der Quellcode für diese Teilaufgabe ist in \lref{2b} gegeben.
\lstinputlisting[firstline=74,lastline=93,label=lst:2b,caption={steuerung.m}]{../code/steuerung.m}

Um zu zeigen, dass der Fehler exponentiell abnimmt, wenn $N$ linear vergrößert
wird, wurde die relative Abweichung vom Mittelwert der 4 Eigenwerte mit den jeweils
größten $N$ geplottet. Das Ergebnis ist in \fref{62b} dargestellt. Man erkennt,
dass der Fehler tatsächlich exponentiell abfällt (im Log-Plot linear) und bei
 $N > 75$ schließlich konstant bleibt. Die untere Grenze für den relativen Fehler
ergibt sich aus der endlichen Maschinengenauigkeit.

\begin{figure}[htb]
  \centering
  \includegraphics[width=0.8\columnwidth,keepaspectratio]{../tmp/plot62b.png}
  \caption{Halblogarthmischer Plot zur Abhängigkeit des Fehlers von der Diskretisierung.}
  \label{fig:62b}
\end{figure}